\summarychapter{35}{A New Order of the Ages}{The Great Awakening}

\begin{inbrief}
In the American colonies the churches entered a new order: no single church could be established, so
each had to win converts and support itself without the state. The Great Awakening, the first of the
American revivals, convinced Christians that this voluntary system could work. The chapter describes
the Puritan background, the revival under Edwards and Whitefield, and the alliance of revivalists and
rationalists that produced religious freedom.
\question{What was the Great Awakening, and why was it crucial to the development of Christianity in
America?}
\end{inbrief}

On July 4, 1776, Franklin, Adams and Jefferson were appointed to design a seal for the United States.
Franklin proposed Moses with Pharaoh's army drowning in the Red Sea and the motto ``Rebellion to
tyrants is obedience to God''; Jefferson, the children of Israel led by cloud and pillar of fire.
Little of this survived, but it is telling that these patriots, themselves distant from their
Jewish-Christian heritage, looked to the Bible for symbols to unite and interpret their people. The
Latin of the seal, still on American currency, means roughly: ``one out of many; God has smiled on
our undertakings; 1776; a new order of the ages.''

\ssection{New Order for the Churches}
For thirty generations Christians had lived by the idea of Christendom: people were born and baptized
into a Christian society in which the church dispensed saving grace and the state kept Christian law
and order. The Reformation shattered it without meaning to, leaving denominations in place of one
church, while national princes continued the alliance of church and state in their own realms, in
churches that swung between repression and laxity. In America the variety of nations and faiths made
an established church implausible in all but a few colonies; by 1646 eighteen languages were heard
along the Hudson alone. Each group wanted liberty to proclaim its own view, and it became plain that
it could have that liberty only by granting it to the rest. The settlers had come to build a
government that would produce holy people; they settled for a society in which people could still
practice holiness. The churches were therefore on their own in evangelizing and nurturing, with no
state support or protection. This condition is voluntaryism. People could accept or reject the
gospel as they pleased, and the denominations had to win converts and raise funds unaided.

The Great Awakening proved that voluntaryism could work, and many came to see revivals as God's gift
for creating a Christian America; Jonathan Edwards even preached that America would be the scene of
the millennium. Resistance was strongest in New England, where Puritanism established by law stood
for the old order. Yet Puritanism carried the seeds of revival. Its church rested on the covenant of
grace and was to be, as far as possible, a church of visible saints, so the first generation required
conversion before membership, a rule reaffirmed in the Cambridge Platform of 1648. The doors of the
church do not stand open to ``all sorts of people good or bad''; applicants must be examined for
repentance and faith and usually gave a personal, public account of God's work in their souls.

\ssection{God's Will in Society}
Massachusetts was more than a string of congregations. Like the Pilgrims in their Mayflower Compact,
the Puritans meant to knit the whole community together by a civil covenant: a Christian people must
walk in God's ways in its government too. God's will was embodied in ``wholesome laws,'' drawn from
Scripture, nature and right reason, which bank the raging sea of human sin; the test of a law was
whether it advanced the public good, that is, God's will in society. Holding the charter, the
Puritans allowed only freemen to vote and required freemen to be church members, so the churches
controlled the vote and public morals. This legislating of morality is one reason later Americans,
who prize personal freedom above the character of society, came to hate them.

The ``holy experiment,'' which joined a church of the truly converted to a Christian state, seemed
bound to fail. As the founders' zeal cooled, fewer people could testify to grace. To keep the
churches from shrinking, many accepted in 1662 the Half-Way Covenant, which gave the ``unawakened'' a
partial membership: they could have their children baptized and share in church life but not take
full Communion. That was enough for most social and political purposes, and the saints shrank to a
small minority. When a new charter in 1691 based the vote on property instead of church membership,
New England stood at a crossroads. Soon two kinds of Puritan heirs appeared. The spiritual heritage
went to the children of the Awakening, with its call for personal conversion. The ``worldly
puritans'' kept the sense of civic responsibility and lawful government and still believed that
empires rise and fall as they obey or disobey Providence, and that God smiles on the quest for
liberty.

For twenty years (1720--1740) the awakening came as regional breezes. Whitefield named Theodore
J.~Frelinghuysen, a Dutch Reformed minister in New Jersey, ``the beginner of the great work'': his
impassioned appeals to the farmers of the Raritan valley brought a large ingathering. The wind passed
to the Scotch-Irish Presbyterians. William Tennent of Pennsylvania trained zealous ministers in a
small school that bookish clergy mocked as the ``Log College,'' and its graduates carried revival
through New Jersey. A quarrel over ``educated'' versus ``converted'' ministers split the Presbyterians
into New Side men, for the revival, and Old Side men, against it. The New Side sent missionaries
south, where Samuel Davies planted revivalist Presbyterian churches in Virginia, and Baptists
multiplied in Virginia and the Carolinas under unschooled, zealous preachers led by Shubael Stearns.
In New England the breeze rose at Northampton, Massachusetts, where Jonathan Edwards had charge of two
hundred families sunk in ``extraordinary dulness in religion.'' In December 1734 ``the Spirit of God
began extraordinarily to set in.'' Scarcely anyone in town was left unconcerned about eternal things,
and souls came ``by flocks to Jesus Christ.''

\ssection{The Great Revivalist}
The winds stayed regional until George Whitefield, the father of modern mass evangelism, arrived in
1739 and preached his way from Georgia to New York. In Philadelphia he spoke out of doors and
impressed even Benjamin Franklin, who marveled that hearers admired a man who assured them they were
``half beasts and half devils,'' and who heard psalms sung in house after house. In 1740 Whitefield
went north and joined the New England revival to that of the Middle Colonies and the South. Boston
ministers invited him, setting the precedent that successful revivalists make for the largest cities.
He spoke in Boston, at Harvard and four times in Edwards's church, then passed through Connecticut,
where people rode in from the farms ``like a steady stream of horses and their riders.'' He left
after a month, and the breezes had become a hurricane, the Great Awakening.

Edwards and other ministers began preaching in neighboring towns. At Enfield, Connecticut, his sermon
``Sinners in the Hands of an Angry God'' was merciless: God holds sinners over the flames as one
holds a loathsome spider over a candle; the ground under them is a rotten floor over a blazing pit.
Sobs and gasps sometimes drowned his voice. The sermon was almost his undoing, for later Americans
remembered him, wrongly, only as a hellfire preacher and forgot the keen psychologist, the brilliant
theologian and the third president of Princeton. By 1741 every element of revival was present:
visiting preachers, warnings of perdition, traveling spellbinders, prayer meetings and a rush of new
members, along with controversies and church splits.

The life of a Connecticut farm boy shows the change. Isaac Backus was taught that good order rests on
the churches' training and the colony's laws. In 1741 the Awakening reached Norwich. His mother was
converted, and Isaac, seventeen, was ``born again'' while mowing alone in a field, without unusual
emotion: by ``divine light'' he saw Christ's perfect righteousness and the riches of his grace. This
``new light'' or inward witness was the key to the New England revival. The revivalists argued that
their fathers had left the Church of England because admitting the unconverted to the church is
contrary to the Word of God, and that the Awakening was a call to a new reformation. New Lights
separated from the parish churches and formed congregations in the founders' way, hearing testimonies
of conversion and signing a covenant as visible saints. Backus became an itinerant evangelist in
southeastern New England, came to Baptist convictions and with a few others formed the First Baptist
Church of Middleborough, Massachusetts.

\ssection{Advocate of Religious Freedom}
More than anyone else Backus formulated and publicized the evangelical view of church and state that
eventually prevailed. In 1769 New England Baptists formed the Warren Association, and two years later
its grievance committee, of which he was the key member, took up the cause of religious freedom. He
wrote tracts and petitions, collected evidence of persecution, testified in court and fought in
newspapers, debates and letters. The imprisonment of his mother, brother and uncle in Connecticut and
his own pietist convictions made him a passionate opponent of establishment. The Baptists held that
every direct tie between the state and institutional religion must be cut so that America might
become truly Christian. Like Jefferson and Madison, Backus believed that ``truth is great and will
prevail,'' but by truth he meant the revealed doctrines of Scripture. God has appointed two kinds of
government, civil and ecclesiastical, ``which ought never to be confounded.'' Legislators are elected
for civil affairs and have no mandate in religion, and in any case ``religion is a voluntary
obedience unto God which force cannot promote.''

The revivalists did not give up the dream of a Christian America. God's kingdom would come if a
majority could be persuaded to submit voluntarily to his laws, and revivals were his means. By 1760
this was no empty hope. From 1740 to 1742 the Awakening had brought 25,000 to 50,000 members into the
New England churches, and between 1750 and 1760 a hundred and fifty new Congregational churches were
formed, besides the multiplying Baptists. Yet the message marked a break with the past. By
concentrating on the individual's need of salvation, the Awakeners neglected the Puritan concern for
the political and social implications of the gospel. With the covenant of grace limited to twice-born
individuals, the idea of a covenant people slid from the church to the American people, and its
mission from a holy commonwealth to the struggle for liberty.

Others favored voluntaryism too. Many Sons of Liberty, such as Jefferson, Franklin and Madison,
shared the Enlightenment's trust in reason and held that a person can be moved only by the weight of
evidence in his own mind; coercion of opinion, said Jefferson, had made ``one half the world fools,
and the other half hypocrites.'' Here the reasonable men and the revival men met, though the
revivalist's freedom was freedom ``in Christ,'' under the Spirit and checked by Scripture. As Sidney
Mead put it, the revivalist ``did not have the heart'' and the rationalist ``did not have the head''
to justify enforced uniformity any longer. Their temporary alliance succeeded in the First
Amendment: ``Congress shall make no law respecting an establishment of religion, or prohibiting the
free exercise thereof.'' A century later Lord James Bryce, in \emph{The American Commonwealth}, noted
that Americans think of the church as a spiritual body with spiritual purposes, so that it never
occurs to them that state churches should exist; such a body wants no state help or privilege.
Nothing in the new order of the ages was newer than that.

\begin{keyfacts}
\keypeople{Jonathan Edwards; George Whitefield; Theodore J.~Frelinghuysen; William Tennent; Samuel
Davies; Shubael Stearns; Isaac Backus; Benjamin Franklin; Thomas Jefferson; James Madison; Lord
Bryce.}
\keydates{%
  \kd{1648}{Cambridge Platform: conversion required for church membership}
  \kd{1662}{Half-Way Covenant}
  \kd{1691}{New Massachusetts charter bases the vote on property}
  \kd{1734}{Revival begins at Northampton under Jonathan Edwards}
  \kd{1739--1740}{Whitefield's preaching tour unites the regional revivals}
  \kd{1741}{Height of the Great Awakening; conversion of Isaac Backus}
  \kd{1769}{New England Baptists form the Warren Association}
  \kd{1776}{American independence; the seal's ``new order of the ages''}}
\keyterms{Christendom; voluntaryism; covenant of grace; civil covenant; visible saints; Half-Way
Covenant; Log College; New Side and Old Side; New Lights; revival; religious freedom; First
Amendment.}
\end{keyfacts}
